\chapter*{T}
\label{chap:t}

\term{T-distribution}
The t-distribution is a probability distribution used when a sample is small and the population variance is unknown. It has heavier tails than the normal distribution, but approaches it as degrees of freedom increase. \par\noindent\textbf{Applications:} it supports small-sample confidence intervals and hypothesis tests.

\term{T-test}
A t-test compares means using an estimate of variability, especially when samples are small. Its assumptions and independent or paired design determine the appropriate version. \par\noindent\textbf{Applications:} t-tests compare treatments, measurements, and experimental groups.

\term{Tableau (Combinatorics)}
A tableau is a filling of a diagram subject to specified ordering rules. Young tableaux use rows and columns to encode partitions and permutations. \par\noindent\textbf{Applications:} tableaux support representation theory, symmetric functions, and combinatorics.

\term{T-norm (Triangular Norm)}
A t-norm is a commutative, associative, monotone operation on $[0,1]$ with $T(a,1)=a$. It generalizes intersection when truth values are graded. \par\noindent\textbf{Applications:} t-norms model conjunction in fuzzy logic and approximate reasoning.

\term{T-conorm}
A t-conorm is a fuzzy-logic operation that generalizes union and is dual to a t-norm, often by $S(a,b)=1-T(1-a,1-b)$. \par\noindent\textbf{Applications:} t-conorms combine graded alternatives in decision systems and control.

\term{Torsion (Group Theory)}
An element of a group is torsion if it has finite order. A torsion group consists entirely of torsion elements. \par\noindent\textbf{Applications:} torsion identifies periodic behavior in algebra, topology, and symmetry groups.

\term{Torsion}
See \term{Torsion Tensor}. \par\noindent\textbf{Applications:} torsion tensors measure the failure of a connection to be symmetric.

\term{Torsion-free Group}
A torsion-free group has no nonidentity element of finite order. The additive group $\mathbb Z$ is an example. \par\noindent\textbf{Applications:} torsion-free groups simplify algebraic and geometric classification.

\term{Torsion Point}
A torsion point on an elliptic curve is a point $P$ for which $nP=\mathcal O$ for some positive integer $n$. \par\noindent\textbf{Applications:} torsion points are important in elliptic-curve arithmetic and cryptography.

\term{Torsion Tensor}
The torsion tensor of a connection is $T(X,Y)=\nabla_XY-\nabla_YX-[X,Y]$. It measures the failure of parallel transport to be symmetric in two directions. \par\noindent\textbf{Applications:} torsion appears in differential geometry, relativity, and theories of defects in materials.

\term{Torsion Subgroup}
The torsion subgroup of an abelian group contains all elements of finite order. It decomposes into primary parts associated with primes. \par\noindent\textbf{Applications:} it separates periodic components from free components in algebraic classification.

\term{Tacnode}
A tacnode is a singular point where two curve branches touch with the same tangent line. \par\noindent\textbf{Applications:} tacnodes arise in algebraic-curve classification, bifurcation diagrams, and geometric modeling.

\term{Tangent}
A tangent describes either the limiting line along a curve at a point or the trigonometric function $\tan x=\sin x/\cos x$. \par\noindent\textbf{Applications:} tangents approximate curves, define slopes, and model angles and periodic systems.

\term{Tangent Space}
The tangent space $T_pM$ is the vector space of directions of curves passing through $p$ on a smooth manifold. \par\noindent\textbf{Applications:} tangent spaces define derivatives, velocities, gradients, and local linear models.

\term{Tangent Bundle}
The tangent bundle $TM$ is the disjoint union of all tangent spaces of a manifold. If $M$ has dimension $n$, then $TM$ has dimension $2n$. \par\noindent\textbf{Applications:} tangent bundles describe vector fields, velocities, and differential equations on manifolds.

\term{Tangent Line}
A tangent line is the best first-order linear approximation to a curve at a point. Its slope is the derivative when the curve is represented as a function. \par\noindent\textbf{Applications:} tangent lines support local approximation, optimization, and motion analysis.

\term{Tangent Vector}
A tangent vector represents the direction and rate of change of a curve at a point. On a smooth manifold it can be defined as a derivation on smooth functions. \par\noindent\textbf{Applications:} tangent vectors represent velocities and infinitesimal changes in geometry and mechanics.

\term{Tautological Bundle}
The tautological bundle over a Grassmannian has fiber $W$ at the point representing the subspace $W$. \par\noindent\textbf{Applications:} it records the subspaces being parametrized and is used in geometry, topology, and characteristic classes.

\term{Tautology (Logic)}
A tautology is a logical formula true under every assignment of its variables. \par\noindent\textbf{Applications:} tautologies support valid inference rules, circuit simplification, and formal verification.

\term{Taylor Series}
A Taylor series expresses a function as an infinite sum of derivative terms at one point, $\sum_{n\ge0}f^{(n)}(a)(x-a)^n/n!$. \par\noindent\textbf{Applications:} Taylor series support approximation, differential equations, numerical methods, and error analysis.

\term{Taylor's Theorem}
Taylor's theorem approximates a function near a point by a polynomial and gives a remainder that bounds the error. \par\noindent\textbf{Applications:} it justifies numerical approximation, local models, and estimates in calculus and differential equations.

\term{Teichmüller Space}
Teichmüller space parametrizes complex structures on a fixed topological surface up to an appropriate equivalence. \par\noindent\textbf{Applications:} it studies Riemann surfaces, moduli spaces, and deformation of geometric structures.

\term{Teichmüller Theory}
Teichmüller theory studies deformations of complex and hyperbolic structures on surfaces. \par\noindent\textbf{Applications:} it connects mapping class groups, Riemann surfaces, hyperbolic geometry, and moduli theory.

\term{T-matrix (Scattering Matrix)}
The T-matrix describes the transition from an incoming state to a scattered state in a scattering problem. \par\noindent\textbf{Applications:} it is used in quantum mechanics, nuclear physics, and wave scattering.

\term{Temperate Distribution}
A tempered distribution is a continuous linear functional on the Schwartz space of rapidly decreasing test functions. It may grow at most polynomially and has a well-defined Fourier transform. \par\noindent\textbf{Applications:} tempered distributions model signals, point sources, and generalized PDE solutions.

\term{Tempered Distribution}
A tempered distribution is a continuous linear functional on the Schwartz space. Examples include the Dirac delta, principal values, and polynomial derivatives. \par\noindent\textbf{Applications:} tempered distributions provide the natural framework for Fourier analysis and weak solutions of PDEs.

\term{Tensor}
A tensor is a multilinear object that generalizes scalars, vectors, and matrices. Its components transform predictably when coordinates change. \par\noindent\textbf{Applications:} tensors describe stress, curvature, material properties, and multilinear data.

\term{Tensor Product}
The tensor product $V\otimes W$ converts bilinear maps from $V\times W$ into linear maps from one vector space. \par\noindent\textbf{Applications:} tensor products combine state spaces, represent multilinear maps, and support quantum mechanics and algebra.

\term{Tensor Algebra}
The tensor algebra of $V$ is the graded algebra $T(V)=\bigoplus_{k\ge0}V^{\otimes k}$, with multiplication given by concatenating tensors. \par\noindent\textbf{Applications:} it constructs free algebras and leads to symmetric and exterior algebras.

\term{Tensor Field}
A tensor field assigns a tensor smoothly to every point of a manifold. \par\noindent\textbf{Applications:} tensor fields represent metrics, electromagnetic fields, stress, and curvature.

\term{Tensor Analysis}
Tensor analysis studies tensors, their coordinate transformations, and their derivatives. \par\noindent\textbf{Applications:} it provides the language of differential geometry, continuum mechanics, and general relativity.

\term{Term (Algebra)}
A term in a polynomial is one coefficient times variables raised to nonnegative integer powers. \par\noindent\textbf{Applications:} terms organize algebraic expressions and support symbolic manipulation and model construction.

\term{Tessellation}
A tessellation partitions a plane or space into cells with no gaps or overlaps. \par\noindent\textbf{Applications:} tessellations model crystals, meshes, spatial data, image processing, and nearest-neighbor geometry.

\term{Ternary Operation}
A ternary operation takes three inputs and returns one output. \par\noindent\textbf{Applications:} ternary operations appear in algebra, logic, programming languages, and multiway composition.

\term{Ternary Ring}
A ternary ring replaces ordinary multiplication by a three-input operation satisfying ring-like axioms. \par\noindent\textbf{Applications:} ternary rings arise in nonassociative algebra and the coordinatization of geometries.

\term{Theorem}
A theorem is a mathematical statement established by a proof from accepted axioms and earlier results. \par\noindent\textbf{Applications:} theorems provide reliable conclusions used in further mathematics and scientific models.

\term{Theory (Logic)}
A theory is a collection of formal statements closed under logical consequence. \par\noindent\textbf{Applications:} theories formalize mathematical structures, axioms, computation, and scientific reasoning.

\term{Theta Function}
A theta function is a holomorphic function with structured periodic or quasi-periodic behavior, commonly defined by a rapidly convergent series. \par\noindent\textbf{Applications:} theta functions occur in elliptic functions, modular forms, abelian varieties, and mathematical physics.

\term{Theta Series}
A theta series is a generating series whose coefficients count representations by a quadratic form, such as sums of squares. \par\noindent\textbf{Applications:} theta series connect number theory, lattices, modular forms, and coding.

\term{Theorema Egregium}
Theorema Egregium states that Gaussian curvature is intrinsic: it can be determined from measurements within the surface and does not depend on how the surface sits in space. \par\noindent\textbf{Applications:} it underlies surface geometry, cartography, and curvature-based shape analysis.

\term{Third Isomorphism Theorem}
The third isomorphism theorem states that if $N\subseteq M$ are normal in $G$, then $(G/N)/(M/N)\cong G/M$. \par\noindent\textbf{Applications:} it simplifies nested quotient constructions in groups, rings, and modules.

\term{Third Root}
The third root, or cube root, of a real number $x$ is the unique real $y$ satisfying $y^3=x$. \par\noindent\textbf{Applications:} cube roots occur in scaling laws, volume calculations, and solving cubic equations.

\term{Thomson's Lamp}
Thomson's lamp is a thought experiment in which a lamp is switched infinitely many times before a finite time. The example shows that an infinite sequence of actions may not determine a final state. \par\noindent\textbf{Applications:} it illustrates limits, supertasks, and questions about infinite processes.

\term{Three-dimensional}
Three-dimensional means that three independent coordinates are needed to specify a position or state. \par\noindent\textbf{Applications:} three-dimensional models describe physical space, graphics, geometry, and engineering structures.

\term{Tilde (Operator)}
A tilde is a notation mark used to indicate a modified, approximate, or related object, such as $\tilde f$. Its meaning depends on context. \par\noindent\textbf{Applications:} tildes denote transforms, estimates, equivalence relations, and asymptotic forms.

\term{Time Complexity}
Time complexity describes how an algorithm's running time grows with input size, often using Big O notation. \par\noindent\textbf{Applications:} it compares algorithms and guides the design of scalable software.

\term{Tikhonov Regularization}
Tikhonov regularization stabilizes an ill-posed inverse problem by minimizing a data-fit term plus a penalty, often $\lambda\|x\|_2^2$. \par\noindent\textbf{Applications:} it reduces noise and overfitting in inverse problems, imaging, and machine learning.

\term{Tiling}
A tiling covers a plane or space with tiles without gaps or overlaps. \par\noindent\textbf{Applications:} tilings appear in architecture, materials, discrete geometry, and computational meshing.

\term{Tiling Theory}
Tiling theory studies how shapes cover spaces and which local rules force or prevent global tilings. \par\noindent\textbf{Applications:} it informs crystallography, aperiodic materials, combinatorics, and pattern design.

\term{Time-reversal Symmetry}
Time-reversal symmetry means that the governing laws remain unchanged when time is replaced by $-t$, with appropriate changes to velocities. \par\noindent\textbf{Applications:} it tests physical models and distinguishes reversible from dissipative dynamics.

\term{Timelike Curve}
A timelike curve in spacetime has a timelike tangent at every point, so it can represent the path of a massive particle. \par\noindent\textbf{Applications:} timelike curves describe causality and possible observer trajectories in relativity.

\term{Timelike Vector}
A timelike vector has negative Lorentzian squared norm under the $(-,+,+,+)$ convention. \par\noindent\textbf{Applications:} timelike vectors represent physical four-velocities and causal directions in relativity.

\term{Tits Building}
A Tits building is a simplicial complex encoding the incidence structure of subgroups associated with a group of Lie type. \par\noindent\textbf{Applications:} buildings connect group theory, geometry, and representation theory.

\term{Tits Alternative}
The Tits alternative states that every finitely generated subgroup of a linear group is either virtually solvable or contains a free group of rank two. \par\noindent\textbf{Applications:} it reveals a major structural divide in linear groups and geometric group theory.

\term{Tits Group}
The Tits group is a finite simple group arising from constructions related to groups of Lie type. \par\noindent\textbf{Applications:} it provides an example in finite-group classification and the theory of sporadic groups.

\term{Toeplitz Matrix}
A Toeplitz matrix is constant along each diagonal from upper left to lower right. \par\noindent\textbf{Applications:} Toeplitz matrices model convolution and support fast algorithms in signal processing and numerical analysis.

\term{Toeplitz Operator}
A Toeplitz operator multiplies a Hardy-space function by a symbol and then projects back to the Hardy space. \par\noindent\textbf{Applications:} Toeplitz operators connect operator theory, signal processing, and boundary-value problems.

\term{Topological Invariant}
A topological invariant is a property unchanged by homeomorphisms, such as connectedness, compactness, or a fundamental group. \par\noindent\textbf{Applications:} invariants distinguish spaces and support classification in topology and geometry.

\term{Topological Space}
A topological space is a set with a collection of subsets called open sets, closed under arbitrary unions and finite intersections and containing the empty set and whole set. \par\noindent\textbf{Applications:} topological spaces formalize continuity, convergence, and geometric structure.

\term{Topology}
Topology studies properties preserved by continuous deformation. It uses open sets to define continuity, convergence, and neighborhoods. \par\noindent\textbf{Applications:} topology supports geometry, data analysis, dynamical systems, and physics.

\term{Topos}
A topos is a category that shares key structural features with the category of sets, including products, function objects, and a classifier for subobjects. \par\noindent\textbf{Applications:} topoi unify geometry and logic and provide settings for sheaves, cohomology, and constructive mathematics.

\term{Total Derivative}
The total derivative accounts for every way a composite function changes when its variables depend on one parameter. \par\noindent\textbf{Applications:} total derivatives support sensitivity analysis, mechanics, optimization, and differential equations.

\term{Total Differential}
The total differential is the linear approximation $df=\sum_i(\partial f/\partial x_i)dx_i$ to a multivariable change. \par\noindent\textbf{Applications:} it estimates measurement error and supports tangent-plane approximations and optimization.

\term{Total Order}
A total order is a partial order in which every two elements are comparable. The usual order on the real numbers is total. \par\noindent\textbf{Applications:} total orders support sorting, optimization, ranking, and induction arguments.

\term{Total Variation}
The total variation of a function is the supremum of sums of absolute changes over all partitions. Finite total variation means the function has controlled cumulative oscillation. \par\noindent\textbf{Applications:} it supports integration, signal analysis, and functions with jumps.

\term{Total Variation Diminishing (TVD)}
Total variation diminishing means that a numerical method does not increase the total variation of the computed solution. \par\noindent\textbf{Applications:} TVD schemes reduce spurious oscillations near shocks in hyperbolic PDEs.

\term{Touchard Polynomials}
Touchard polynomials encode the numbers of partitions of a set into blocks, using Stirling numbers as coefficients. \par\noindent\textbf{Applications:} they support combinatorics, moments of Poisson variables, and enumerative probability.

\term{Tournament (Graph Theory)}
A tournament is a directed graph obtained by assigning one direction to every edge of a complete graph. \par\noindent\textbf{Applications:} tournaments model competitions, preference comparisons, and ranking systems.

\term{Torsion-free Abelian Group}
A torsion-free abelian group has no nonidentity element of finite order. \par\noindent\textbf{Applications:} torsion-free groups describe lattice-like structures and the free part of abelian-group classifications.

\term{Trace}
The trace of a square matrix is the sum of its diagonal entries. It is invariant under similarity and equals the sum of eigenvalues with multiplicity. \par\noindent\textbf{Applications:} trace is used in linear algebra, covariance analysis, quantum mechanics, and optimization.

\term{Transcendence Degree}
The transcendence degree of $L/K$ is the size of a largest algebraically independent subset of $L$ over $K$. It measures the number of independent transcendental parameters. \par\noindent\textbf{Applications:} it gives the dimension of algebraic varieties and supports field theory and transcendence results.

\term{Transcendental Extension}
A transcendental extension contains an element that does not satisfy any nonzero polynomial with coefficients in the base field. \par\noindent\textbf{Applications:} transcendental extensions describe function fields and dimensions in algebraic geometry.

\term{Transfinite Induction}
Transfinite induction proves statements on a well-ordered set by proving each case from all earlier cases. \par\noindent\textbf{Applications:} it supports proofs about ordinals, cardinals, recursive constructions, and set theory.

\term{Transfinite Number}
A transfinite number describes the size or order type of an infinite set, using cardinals or ordinals such as $\aleph_0$. \par\noindent\textbf{Applications:} transfinite numbers organize infinite hierarchies in set theory and logic.

\term{Transfinite Recursion}
Transfinite recursion defines an object at each ordinal using values already defined at smaller ordinals. \par\noindent\textbf{Applications:} it constructs functions and hierarchies in set theory, logic, and infinitary combinatorics.

\term{Transformation}
A transformation is a function that maps one mathematical object or space to another. \par\noindent\textbf{Applications:} transformations describe coordinate changes, geometric motions, data processing, and symmetries.

\term{Transformation (Linear)}
A linear transformation preserves vector addition and scalar multiplication. \par\noindent\textbf{Applications:} linear transformations represent matrices, physical systems, projections, and changes of coordinates.

\term{Transformation (Möbius)}
A Möbius transformation has the form $f(z)=(az+b)/(cz+d)$ with $ad-bc\ne0$. It acts conformally on the extended complex plane. \par\noindent\textbf{Applications:} Möbius maps support complex analysis, hyperbolic geometry, and fractional-linear models.

\term{Transitive Action}
An action of $G$ on $X$ is transitive when one group element can carry any point to any other point. \par\noindent\textbf{Applications:} transitive actions describe homogeneous spaces, symmetries, and geometric configurations.

\term{Translation}
A translation moves every point by the same vector, preserving distances and directions. \par\noindent\textbf{Applications:} translations model displacement, coordinate changes, image registration, and periodic structures.

\term{Transpose}
The transpose $A^T$ is obtained by exchanging rows and columns, so $(A^T)_{ij}=A_{ji}$. \par\noindent\textbf{Applications:} transposes define symmetry, inner products, least squares, and matrix factorizations.

\term{Transversality}
Two submanifolds are transverse when their tangent spaces together span the tangent space of the ambient manifold at each intersection point. \par\noindent\textbf{Applications:} transversality makes intersections stable and supports differential topology and geometry.

\term{Transverse Intersection}
A transverse intersection satisfies the transversality condition, so the tangent spaces meet in the expected dimension. \par\noindent\textbf{Applications:} transverse intersections are stable under perturbation and are used to define intersection numbers.

\term{Transposition (Permutation)}
A transposition is a permutation that swaps two elements and fixes all others. Every finite permutation is a product of transpositions. \par\noindent\textbf{Applications:} transpositions generate symmetric groups and support permutation algorithms.

\term{Transport (Parallel)}
Parallel transport moves a vector along a curve while keeping its covariant derivative zero. \par\noindent\textbf{Applications:} it compares tangent vectors at different points and defines holonomy, curvature, and geometric phase.

\term{Trapezium}
A trapezium is a quadrilateral whose definition varies by convention: it may have no parallel sides or at least one pair. \par\noindent\textbf{Applications:} trapezia are used in geometry, area calculation, and design.

\term{Trapezoidal Rule}
The trapezoidal rule approximates an integral by replacing each interval with a trapezoid and summing the areas. \par\noindent\textbf{Applications:} it evaluates experimental data and numerical integrals when antiderivatives are unavailable.

\term{Tree}
A tree is a connected graph with no cycles. A tree on $n$ vertices has $n-1$ edges. \par\noindent\textbf{Applications:} trees model hierarchies, file systems, search algorithms, and communication networks.

\term{Triangle Inequality}
The triangle inequality states that $d(x,z)\leq d(x,y)+d(y,z)$ for a metric. A direct route is no longer than a route through an intermediate point. \par\noindent\textbf{Applications:} it controls convergence, error bounds, clustering, and shortest-path algorithms.

\term{Triangulated Category}
A triangulated category is an additive category with a shift operation and distinguished triangles satisfying axioms modeled on mapping cones. \par\noindent\textbf{Applications:} triangulated categories organize derived categories, homological algebra, algebraic geometry, and stable homotopy.

\term{Trivial Group}
The trivial group contains only the identity element and has order one. \par\noindent\textbf{Applications:} it is the base case for group extensions, kernels, and induction arguments.

\term{Trivial Ring}
The trivial ring has one element, so $0=1$. \par\noindent\textbf{Applications:} it serves as a zero object or degenerate case in ring theory and algebraic geometry.

\term{Trivial Topology}
The trivial topology contains only the empty set and the whole space as open sets. \par\noindent\textbf{Applications:} it provides simple examples and counterexamples in topology and continuity.

\term{Truncation Error}
Truncation error is the error introduced when an exact equation or infinite process is replaced by a finite approximation. \par\noindent\textbf{Applications:} it guides step-size selection and accuracy analysis in numerical methods.

\term{Turán's Theorem}
Turán's theorem gives the maximum number of edges in an $n$-vertex graph that contains no complete graph on $r+1$ vertices. \par\noindent\textbf{Applications:} it bounds dense networks and is foundational in extremal graph theory.

\term{Turing Machine}
A Turing machine is an abstract model with a tape, read-write head, finite control, and transition rules. It captures the idea of algorithmic computation and has undecidable problems such as halting. \par\noindent\textbf{Applications:} Turing machines define computability and complexity classes.

\term{Twin Primes}
Twin primes are pairs of primes differing by two, such as $(3,5)$. The twin-prime conjecture says that infinitely many such pairs exist. \par\noindent\textbf{Applications:} they are a central example in prime-distribution research.

\term{Two Envelopes Paradox}
The two-envelopes paradox appears to show that switching envelopes always increases expected money, even though one envelope is already selected. The error comes from using an unjustified probability model for unbounded amounts. \par\noindent\textbf{Applications:} it clarifies conditional expectation and decision-making under uncertainty.

\term{Two's Complement}
Two's complement represents signed binary integers so that negation is performed by inverting bits and adding one. \par\noindent\textbf{Applications:} it allows computer hardware to use the same addition circuit for positive and negative integers.

\term{Tychonoff's Theorem}
Tychonoff's theorem states that any product of compact spaces is compact in the product topology. \par\noindent\textbf{Applications:} it supports infinite-dimensional analysis, functional analysis, and existence proofs.

\term{Type Theory}
Type theory assigns types to expressions to organize computation and proofs. Under the Curry--Howard correspondence, types represent propositions and programs represent proofs. \par\noindent\textbf{Applications:} type theory underlies programming languages, proof assistants, formal verification, and foundations of mathematics.
\term{Stirling Number}
Stirling numbers of the first kind count permutations by cycles, while those of the second kind count partitions into nonempty blocks. \par\noindent\textbf{Applications:} they expand factorial polynomials and support enumerative combinatorics and probability.
\term{Strong Duality}
Strong duality means that a primal optimization problem and its dual have equal optimal values. \par\noindent\textbf{Applications:} it provides certificates of optimality and supports algorithms in linear, convex, and machine-learning optimization.
